\documentclass[12pt,a4paper,UTF8]{ctexart}
\usepackage[a4paper,top=22mm,bottom=22mm,left=22mm,right=22mm,headheight=15pt]{geometry}
\usepackage{amsmath,amssymb,mathtools,bm}
\usepackage{graphicx,xcolor,booktabs,longtable,array,enumitem,fancyhdr,caption,float,microtype,setspace}
\usepackage[hidelinks,unicode]{hyperref}
\usepackage[most]{tcolorbox}
\definecolor{ContestBlue}{HTML}{17365D}
\hypersetup{pdftitle={第087套},pdfauthor={}}
\setstretch{1.24}
\setlength{\parindent}{2em}
\setlength{\parskip}{0.35em}
\setlength{\emergencystretch}{2em}
\setlist{itemsep=0.35em,topsep=0.4em}
\captionsetup{font=small,labelfont=bf,justification=centering}
\setcounter{secnumdepth}{0}
\pagestyle{fancy}
\fancyhf{}
\fancyhead[L]{\small 第087套}
\fancyfoot[C]{\thepage}
\renewcommand{\headrulewidth}{0.4pt}
\ctexset{
  section={format=\Large\bfseries\color{ContestBlue},beforeskip=1.2em,afterskip=0.7em},
  subsection={format=\large\bfseries\color{ContestBlue},beforeskip=1.0em,afterskip=0.5em},
  subsubsection={format=\normalsize\bfseries,beforeskip=0.8em,afterskip=0.35em}
}
\makeatletter
\newcommand{\tightlist}{\setlength{\itemsep}{0pt}\setlength{\parskip}{0pt}}
\newcommand{\maxwidth}{\ifdim\Gin@nat@width>\linewidth\linewidth\else\Gin@nat@width\fi}
\newcommand{\maxheight}{\ifdim\Gin@nat@height>0.72\textheight 0.72\textheight\else\Gin@nat@height\fi}
\newcommand{\pandocbounded}[1]{\begingroup\setkeys{Gin}{width=\maxwidth,height=\maxheight,keepaspectratio}#1\endgroup}
\makeatother
\begin{document}
\begin{titlepage}
  \centering
  \vspace*{58mm}
  {\Huge\bfseries\color{ContestBlue} 第087套\par}
  \vfill
  \vspace*{28mm}
\end{titlepage}
\section{第 1 题}\label{ux7b2c-1-ux9898}

两个相同小球的质量均为 \(m\)，用一根质量不计、不可伸长、长度为 \(l\)
的细线连接，放在光滑的水平面上，如图所示。给右侧小球一个大小为
\(v_0\)、方向竖直向上的初速度，重力加速度为 \(g\)。

\begin{enumerate}
\def\labelenumi{\arabic{enumi}.}
\tightlist
\item
  当细线方向变为竖直时，求上方小球轨迹在该时刻的曲率半径；
\item
  若要求在这一时刻下方小球对水平面的压力恰好为零，求初速度 \(v_0\)。
\end{enumerate}

\begin{figure}[H]
\centering
\includegraphics[width=0.9\linewidth,height=\textheight,keepaspectratio,alt={图 1-1　两小球、细线及初速度的初始位置。}]{../figures/question-01.pdf}
\caption*{图 1-1　两小球、细线及初速度的初始位置。}
\end{figure}

\newpage

\section{第 2 题}\label{ux7b2c-2-ux9898}

两个圆的圆心分别位于 \((\pm r,0)\)，半径均为

\[
a=\frac{2\sqrt3}{3}r.
\]

两圆区域内存在方向相反、大小均为

\[
B=B_0\frac{t}{t_0}
\]

的匀强磁场，方向如图所示。

\begin{enumerate}
\def\labelenumi{\arabic{enumi}.}
\tightlist
\item
  求两圆相交部分的感生电场分布；
\item
  在图示电路中，各电压表的上端接红表笔、下端接黑表笔，电压表内阻均为
  \(2R_0\)。求三个电压表的读数以及 \(U_a-U_b\)。
\end{enumerate}

\begin{figure}[H]
\centering
\includegraphics[width=0.9\linewidth,height=\textheight,keepaspectratio,alt={图 2-1　两圆形时变磁场与电压测量电路。}]{../figures/question-02.pdf}
\caption*{图 2-1　两圆形时变磁场与电压测量电路。}
\end{figure}

\newpage

\section{第 3 题}\label{ux7b2c-3-ux9898}

金刚石晶体可由下述方式理解：在立方晶胞的 \(8\) 个顶点和 \(6\)
个面心放置碳原子；相邻三个面的面心与它们最近的顶点构成一个正四面体，共可选出
\(8\) 个这样的正四面体。在其中两两不相邻的 \(4\)
个正四面体中心各放置一个碳原子，再把相邻碳原子用碳碳键连接，即得到金刚石结构。

现在把每个碳原子替换成电路节点，把每条碳碳键替换成阻值为 \(r\)
的电阻。在无限大的金刚石电阻网络中，求：

\begin{enumerate}
\def\labelenumi{\arabic{enumi}.}
\tightlist
\item
  相邻两个节点之间的等效电阻；
\item
  次相邻两个节点之间的等效电阻。
\end{enumerate}

\begin{figure}[H]
\centering
\includegraphics[width=0.517\linewidth,height=\textheight,keepaspectratio,alt={图 3-1　金刚石晶格的三维四面体连接。}]{../figures/question-03.pdf}
\caption*{图 3-1　金刚石晶格的三维四面体连接。}
\end{figure}

\newpage

\section{第 4 题}\label{ux7b2c-4-ux9898}

水平桌面上有两条光滑抛物线导轨，满足

\[
y=kx^2,
\]

其中 \(k\) 为常量，\(xOy\) 平面与桌面重合。一根平行于 \(x\)
轴的导体棒横跨两导轨，棒的两端始终与导轨接触。空间中存在竖直方向、磁感应强度为
\(B\) 的匀强磁场。导体棒从坐标原点开始沿 \(y\) 轴正方向以恒定速度 \(v\)
运动。

\begin{enumerate}
\def\labelenumi{\arabic{enumi}.}
\tightlist
\item
  求导体棒上的动生电动势 \(\mathcal E(t)\)；
\item
  若导体棒单位长度电阻为 \(\lambda\)，导轨电阻不计，求回路电流；
\item
  求维持导体棒匀速运动所需外力的功率。
\end{enumerate}

\begin{figure}[H]
\centering
\includegraphics[width=0.494\linewidth,height=\textheight,keepaspectratio,alt={图 4-1　抛物线导轨、导体棒及磁场方向。}]{../figures/question-04.pdf}
\caption*{图 4-1　抛物线导轨、导体棒及磁场方向。}
\end{figure}

\newpage

\section{第 5 题}\label{ux7b2c-5-ux9898}

一只企鹅在距离水平冰面高度为 \(h_0\) 的位置，以初速度 \(v_0\)
喷出粪便。随后采用下列流体模型研究其喷射过程。

\begin{enumerate}
\def\labelenumi{\arabic{enumi}.}
\tightlist
\item
  忽略空气阻力，求使水平射程最大的喷射仰角
  \(\alpha\)。以水平面上方为正角、下方为负角；
\item
  把粪便视为密度为
  \(\rho\)、不可压缩且无黏性的液体。企鹅利用胃内高压，使充满整个直肠的液体由泄殖孔喷出，出口速度为
  \(v_0\)。直肠近似为长度 \(l\)、半径 \(r\)
  的光滑均匀圆管，胃与泄殖孔的高度差为 \(h\)，大气压为
  \(P_0\)，初始胃内压强为 \(P_0+P_1\)，求 \(P_1\)；
\item
  若考虑黏滞，黏滞系数为 \(\eta\)，直肠中形成充分发展的层流，满足
\end{enumerate}

\[
\Delta P=\frac{8\eta lQ}{\pi r^4},
\]

且速度随距中心轴的距离二次下降、在管壁处为零。求中心轴速度为 \(v_0\)
时的体积流率 \(Q_1\)； 4. 在第（3）问模型下，若中心轴速度仍为
\(v_0\)，初始胃内压强为 \(P_0+P_2\)，求 \(P_2\)； 5.
回到第（2）问的无黏、不可压缩模型。把直肠两端封闭，在体积为 \(V\)
的胃中充入压强 \(P_0\)、温度 \(T_0\)
的双原子理想气体。直肠中靠近胃的一端装入长度为 \(l/2\)
的绝热粪便，忽略直肠高度 \(h\)。迅速压缩胃内气体，使其体积缩小到
\(nV\)，随后用该气体把粪便推出。求粪便喷出时的速度，结果以
\(\sqrt{P_0/\rho}\) 为单位。已知 \(\pi r^2l<(1-n)V\)。

参数为

\[
\begin{gathered}
v_0=2.0\,\mathrm{m/s},\quad g=9.8\,\mathrm{m/s^2},\quad h_0=0.192\,\mathrm m,\\
l=0.20\,\mathrm m,\quad h=0.18\,\mathrm m,\quad r=1.0\times10^{-3}\,\mathrm m,\\
\eta=4.04\times10^{-3}\,\mathrm{Pa\cdot s},\quad
\rho=922\,\mathrm{kg/m^3},\\
\frac{V}{\pi r^2l}=129,\qquad n=0.587.
\end{gathered}
\]

\begin{figure}[H]
\centering
\includegraphics[width=0.43\linewidth,height=\textheight,keepaspectratio,alt={图 5-1　喷射抛体与直肠流动模型。}]{../figures/question-05.pdf}
\caption*{图 5-1　喷射抛体与直肠流动模型。}
\end{figure}

\newpage

\section{第 6 题}\label{ux7b2c-6-ux9898}

日落时，太阳光实际从当地水平线以下射来，但由于大气折射，地面观察者仍能看见太阳。设某行星半径为
\(R_0\)，大气球对称且十分厚重，折射率与离地高度 \(h\) 的关系为

\[
n(h)=\sqrt{5-\frac{(h+R_0)^2}{R_0^2}},
\qquad 0<h<R_0.
\]

当地面观察者刚好看见太阳落山时，求太阳的实际方向与观察者所认为的太阳方向之间的夹角
\(\beta\)。

\begin{figure}[H]
\centering
\includegraphics[width=0.502\linewidth,height=\textheight,keepaspectratio,alt={图 6-1　球对称分层大气中的日落光线。}]{../figures/question-06.pdf}
\caption*{图 6-1　球对称分层大气中的日落光线。}
\end{figure}

\newpage

\section{第 7 题}\label{ux7b2c-7-ux9898}

冥王星半径为 \(r_1=2.30\times10^6\,\mathrm m\)，质量为
\(m_1=1.29\times10^{22}\,\mathrm{kg}\)，自转周期为 \(T=6.387\)
天。卡戎是冥王星的伴星，距离冥王星中心
\(l=1.974\times10^7\,\mathrm m\)，半径为
\(r_2=6.03\times10^5\,\mathrm m\)。假设卡戎与冥王星做圆周运动，公转周期等于冥王星自转周期。取

\[
G=6.67\times10^{-11}\,\mathrm{N\,m^2/kg^2}.
\]

\begin{enumerate}
\def\labelenumi{\arabic{enumi}.}
\tightlist
\item
  求卡戎质量 \(m_2\)，保留两位有效数字；
\item
  设有一条连接两星体的``天梯''，把冥王星上的少量质量
  \(\Delta m=1.0\times10^{10}\,\mathrm{kg}\)
  缓慢搬运到卡戎。求系统周期的增加量 \(\Delta T\)
  和两星距离随轨道角度的平均值变化量
  \(\Delta l\)，均保留两位有效数字。搬运后轨道不再严格为圆轨道。
\end{enumerate}

\begin{figure}[H]
\centering
\includegraphics[width=0.661\linewidth,height=\textheight,keepaspectratio,alt={图 7-1　冥王星---卡戎同步圆轨道与质量输运。}]{../figures/question-07.pdf}
\caption*{图 7-1　冥王星---卡戎同步圆轨道与质量输运。}
\end{figure}

\newpage

\section{第 8 题}\label{ux7b2c-8-ux9898}

黑体吸收所有射向其表面的电磁波，也不断向外辐射电磁波。单位时间、单位面积的黑体辐射功率为

\[
q=\sigma T^4.
\]

若物体对入射电磁波的反射比例为 \(r\)，由基尔霍夫定律，其本征辐射功率变为

\[
q'=(1-r)\sigma T^4.
\]

平面向外辐射时，在与法线夹角为 \(\theta\)
的方向上，单位立体角辐射强度正比于 \(\cos\theta\)。

现有一个温度为 \(T\)、比热容为 \(C\)、质量为 \(m\)
的薄圆盘。一面为黑体，另一面的反射系数为
\(r\)。把圆盘置于真空中，它因两面辐射不对称而获得动量。忽略相对论效应和宇宙背景辐射，求圆盘最终能够获得的最大动能。光子能量和动量满足

\[
E=h\nu,
\qquad
p=\frac{h\nu}{c}.
\]

\begin{figure}[H]
\centering
\includegraphics[width=0.616\linewidth,height=\textheight,keepaspectratio,alt={图 8-1　两面发射率不同的高温圆盘。}]{../figures/question-08.pdf}
\caption*{图 8-1　两面发射率不同的高温圆盘。}
\end{figure}
\end{document}
